\section{Background and Threat Model}
\label{sec:threat}

\subsection{Setting}
A user $u$ delegates an analytics task, stated in natural language, to an LLM agent. The agent answers by issuing
SQL through a query tool. Without our mechanism, the agent's database credential carries $u$'s full
entitlement $E(u)$ (condition B0 in our evaluation): it can read every column and row $u$ can read, regardless of
what the task needs. The trusted broker that provisions the agent also sees the task text and the schema.

\subsection{Adversary}
The adversary wants the agent to read, and reveal, data that the current task does not need. We consider an
adversary who can:
\begin{enumerate}
\item place instructions in content the agent reads: database text fields returned by queries, and documents or
notes attached to the task context (indirect prompt injection~\cite{greshake2023not});
\item thereby make the agent issue arbitrary SQL on its connection, including attempts to bypass enforcement
(catalog reads, role switching, \texttt{COPY}, leaky functions, error and timing side channels);
\item make the agent file widening requests with plausible reasons;
\item observe the agent's final answer (the exfiltration channel we measure).
\end{enumerate}
The adversary cannot modify the broker, the database configuration or the user's entitlement, and holds no
database credential of its own.

\subsection{Trust assumptions}
We rely on A1--A4 (Section~\ref{sec:assumptions}): the DBMS enforces grants, row policies, view ownership and
\texttt{security\_barrier} correctly; the ephemeral login is the agent's only credential; the broker is trusted;
and no involved role has bypass privileges.

\subsection{In scope and out of scope}
\textbf{In scope:} reducing the data an agent can \emph{reach} for a task, measured as over-exposure; containing
injected requests for out-of-task data; bounding and auditing widening; and robustness of the enforcement
against the bypass attempts listed above.

\textbf{Out of scope (not protected):}
(i)~inference and aggregation attacks that use data inside the task scope or inside $E(u)$~\cite{denning1979tracker,adam1989security};
(ii)~covert and timing channels;
(iii)~exfiltration of data that is inside the task scope (TSE narrows what can leak; it does not stop the agent
from repeating in-scope data);
(iv)~a malicious authorized user, who already holds $E(u)$;
(v)~a compromised broker or a misconfigured entitlement $E(u)$;
(vi)~schema metadata: table and column names are visible in the PostgreSQL catalog to every role.
TSE therefore reduces the data reachable per task; it does not prevent all leakage.
